% year: תשע"ח
% type: מבחן
% semester: ב
% moed: א
% number: 1ג
% points: 5
% categories: func-analysis, derivatives
% source: exams/מבחנים/תשעח/מועד א תשעח סמסטר ב.pdf

\begin{question}
(5 נק') עבור איזה ערכים של $a$ לגרף הפונקציה $f(x)=x^4-ax^2$ אין נקודות פיתול?
\end{question}

\begin{hint}
נקודת פיתול היא נקודה שבה $f''$ מחליפה סימן. חשבו את $f''$ ובדקו לאילו ערכים של $a$ היא מחליפה סימן.
\end{hint}

\begin{solution}
$f''(x)=12x^2-2a$. נקודת פיתול היא נקודה שבה $f''$ מחליפה סימן.
\begin{itemize}
  \item אם $a>0$: $f''(x)=0\iff x=\pm\sqrt{a/6}$, ו-$f''$ (פרבולה) מחליפה סימן בכל אחת מהנקודות האלה, לכן יש שתי נקודות פיתול.
  \item אם $a=0$: $f''(x)=12x^2\ge0$, מתאפסת רק ב-$x=0$ בלי להחליף סימן — אין נקודות פיתול ($f=x^4$ קמורה).
  \item אם $a<0$: $f''(x)=12x^2-2a>0$ לכל $x$ — אין נקודות פיתול.
\end{itemize}
\textbf{תשובה:} $a\le0$.
\end{solution}
